
\begin{resumo}[ABSTRACT]
\begin{otherlanguage}{english}
\vspace*{-6mm}

Operating a process plant involves more than tuning loops: one has to choose the operating point, decide what to keep constant when disturbances arrive, control the plant and, finally, check whether the plant does what was decided. This work deals with this last step, supervision, and uses the term operating policy for the set of decisions the plant must follow. The goal is to show that a supervisory layer built on Kubernetes can follow a plant from a policy declared by the user, without the supervisor knowing the plant, using the Tennessee Eastman Process (TEP) as the plant and OPC~UA as the communication standard. Since the work belongs to the Control System Integration track, the focus is on connecting the layers, not on tuning controllers. The Downs and Vogel model was rewritten in Rust as a plant that runs continuously and publishes its signals over OPC~UA. A historian written in Python collects these signals and computes window statistics and a control loop performance index. A custom controller, written in Go and running on a local Kubernetes cluster, reads three resources declared by the user (the cost function, the operating policy and the plant), evaluates the plant against the Mode~1 policy every 30~s and writes the verdict into Kubernetes itself, where a web HMI reads it. The plant reproduced the Downs and Vogel base case, with a cost of 170.63~\$/h against the published 170.61~\$/h, and, in the version with integrated cooling water, Ricker's Mode~1 optimum, with 114.32~\$/h against 114.31~\$/h. With the policy fixed, the verdict stayed compliant in normal operation and flagged the degradation caused by disturbance IDV(6) 1.15~h after it, in the physical order of events: first the product flow, then the cost. The loop index alone did not notice the disturbance, and a setpoint tracking criterion had to be added. A re-reading of the data showed that the index was computed on an inverted scale; on the original article's scale, it would have flagged the disturbance before the economic verdict, but it also flags harmless transients. In a third experiment, with PI loops and no reactor level loop, the plant was lost in about 33~h while the cost went down, and only the level limit declared in the policy noticed it, and late. The objective was reached, with limits: the desired-state and observed-state mechanism of Kubernetes proved suitable for supervision, but a useful policy must also watch inventories and trends, not only cost and loops. Making the supervisor act on the plant is left for future work.

\textbf{Keywords}: Control system integration. Tennessee Eastman Process. Process supervision. Operating policy. Kubernetes. OPC UA.
  \end{otherlanguage}
\end{resumo}
